%! Sinusoidal Functions — Unit 6, Lesson 6.6 — Homework
\documentclass[10pt]{article}
\usepackage{precalculus-article}
\usepackage{precalculus-boxes}
\newcommand{\TallMath}[1]{$\displaystyle #1\rule[-1.4em]{0pt}{3.2em}$}

\begin{document}

\pageheader{Unit 6, Lesson 6.6}{Homework}
\namedateperiod

\vspace{0.12in}

\begin{notesbox}{Sinusoidal Functions --- Practice}

\textbf{1.} A sinusoidal function has a maximum value of $10$ and a minimum value
of $-2$.

\medskip
(a) Find the amplitude. \writelines{2}

(b) Find the midline.  Write your answer as an equation $y = \ldots$ \writelines{2}

\vspace{0.08in}

\textbf{2.} The period of a sound wave is $0.002$ seconds.

\medskip
(a) Find the frequency of the wave in cycles per second (Hz). \writelines{2}

(b) A second wave has frequency $250$ Hz.  What is its period? \writelines{2}

\vspace{0.08in}

\textbf{3.} Is $f(\theta) = \sin(2\theta)$ an odd function, an even function, or
neither?  Justify your answer by computing $f(-\theta)$ and comparing it to
$f(\theta)$ and $-f(\theta)$.

\writelines{4}

\vspace{0.08in}

\textbf{4.} \textit{AP-Style Multiple Choice.}

Which of the following is the midline of a sinusoidal function with maximum
value $6$ and minimum value $-2$?

\medskip
\begin{tabular}{@{}lllll@{}}
(A) $y = 4$ \quad & (B) $y = 2$ \quad & (C) $y = 3$ \quad & (D) $y = -2$ \quad & (E) $y = 6$
\end{tabular}

\vspace{0.08in}
Your answer: \blank{1.5cm} \quad Show your work:

\writelines{3}

\vspace{0.08in}

\textbf{5.} The frequency of a sinusoidal function is $\dfrac{3}{2\pi}$.
What is the period of this function?

\writelines{3}

\vspace{0.08in}

\textbf{6.} A sinusoidal function has amplitude $5$ and midline $y = -3$.
What are the maximum and minimum values of the function?

\writelines{3}

\end{notesbox}

\vspace{0.08in}

\begin{extensionbox}
\textbf{Extension:} A sinusoidal function is simultaneously an odd function
\textit{and} an even function --- that is, $f(-\theta) = -f(\theta)$ and
$f(-\theta) = f(\theta)$ both hold for all $\theta$.  What must be true about
the amplitude and all output values of this function?

\writelines{3}
\end{extensionbox}

\vspace{0.08in}

\begin{tcolorbox}[colback=goldbg, colframe=goldacc, arc=2mm,
    left=3mm, right=3mm, top=2mm, bottom=2mm]
\textbf{Preview of Lesson 3.6 --- Sinusoidal Function Transformations}

In the next lesson you will explore the general form
$g(\theta) = a\sin(b(\theta + c)) + d$.
The parameters $a$, $b$, $c$, and $d$ each control one of the characteristics
you studied today. As you finish this homework, think about which parameter
might control amplitude, which controls period, and which determines the midline.
\end{tcolorbox}

\end{document}
